%% =====================================================================
%%  T-10-02  The Reactance Trigger
%%  -------------------------------------------------------------------
%%  One idea per file. Core is mandatory. Slots in canonical order.
%%  Max 700 words. No layout commands. No filler. New source material
%%  for this entry goes in sources/<BOOK>/T-10-02.tex -- never by
%%  rewriting this file (docs/RULES.md R0.1).
%% =====================================================================
%% @kind: topic
%% @id: T-10-02
%% @title: The Reactance Trigger
%% @subtitle: A threatened freedom is defended out of proportion to its value
%% @chapter: c10
%% @order: 20
%% @status: stable
%% @sources: B4
%% @updated: 2026-09-19

\topic{T-10-02}{The Reactance Trigger}{A threatened freedom is defended out of proportion to its value}

\begin{core}
When a freedom is limited, people want it more than they wanted it before, and want whatever it gave them access to more than they valued it before. The reaction is against the restriction, not in favour of the thing --- which is why it produces value that was never there.
\end{core}
\begin{mechanism}
The mechanism is a defence of the range of options rather than of any particular option. A person who loses an ability loses part of their control over their own circumstances, and control is defended disproportionately because it underwrites everything else.

Three properties follow, and they are what make the trigger usable. First, the escalation is real but displaced: the thing becomes more desired because it is the site of the lost freedom, and the desire survives the loss. Second, the trigger is most powerful where the freedom was already held and is then removed --- a restriction on something never permitted produces far less than a restriction on something previously ordinary. Third, the reaction is intensified by competition: scarcity that is also contested combines the loss of an option with the presence of a rival, and both fire at once.

The practical consequence is that withdrawal works better than refusal. Take something away and it is defended; never offer it and it is not missed.
\end{mechanism}
\begin{tells}
  \item You want the thing more since it became harder to get.
  \item The restriction arrived after the freedom was established.
  \item Others are competing for the same limited thing.
  \item Your objection is to being restricted rather than to the substance.
  \item You cannot say what you would do with the thing if you had it.
\end{tells}
\begin{moves}
  \item Grant a freedom and then withdraw it, rather than refusing it from the start.
  \item Withdraw gradually, so each stage is felt as a new loss.
  \item Introduce competition for what remains.
  \item Frame the restriction as imposed by circumstance rather than chosen by you --- an externally caused loss is still defended, and the defender does not turn on you.
\end{moves}
\begin{counters}
  \item Ask what the thing was worth before it was restricted. That is still its worth; the increase belongs to the restriction.
  \item Identify which freedom you are actually defending, and check whether the thing on the table is what that freedom was for.
  \item Notice competition and discount it. A rival increases the price without increasing the value.
  \item Test by accepting the restriction. If the desire falls away rather than intensifying once you stop resisting, it was reactance and not preference.
  \item Ask who benefits from your defending this particular freedom.
\end{counters}
\begin{cost}
Deliberately withdrawing a freedom to generate desire is the point at which this stops being persuasion, and the counterpart usually knows it --- the feeling of being managed is distinctive and is rarely forgiven. It also produces compliance of a particular bad kind: the person obtains the thing and immediately loses interest, because what they wanted was the freedom rather than the object. That is a poor outcome for a seller, a worse one for a partner, and worst of all when turned inward, where a person who restricts themselves to manufacture motivation ends up wanting nothing they can actually have.
\end{cost}

\sources{B4}
\seesources
\seealso{T-10-01, T-10-03, T-24-03, T-21-01}
